\scanpage{275}\mspage{191}
% archive sheet number: 274
et que je vérifie si cette erreur \unsure{entache|entâche}
\deleted{\unsure{entache d'erreur}} rend inopérantes certaines
démonstrations, voire même, rend faux certains
énoncés. \unsure{Aussi|Ainsi}, il me faudrait donner
un exemple où $C$ est str. gén., mais
$C'\supset C$ ne l'est pas.

Finalement, réflexion faite, ces ennuis \unsure{craints|criants}
disparaissent. Je vais énoncer un résultat
en forme :

\deleted{Lemme 4.8.1}
\uline{Proposition} 4.8.2. Soient $C,M$ deux \inserted{$\unsure{U|\mathcal U}$-}catégories,
$i:C\unsure{\longrightarrow|\Longrightarrow} M$ un foncteur, d'où un
foncteur
\[
\left|\begin{aligned}
\varphi &: M\longrightarrow C^\wedge\\
\varphi(\unsure{X|x})&=(x\longmapsto\operatorname{Hom}_M(i(x),X)
\end{aligned}\right.
\tag*{(4.8.2.1)}
\]
(foncteur composé $M\xrightarrow{\mathrm{can}}M^\wedge
\xrightarrow[\text{i.e. composition avec }i]{i^*}C^\wedge$).

\deleted{$1^{\mathrm{o}}$) Considérons les conditions suivantes}
Pour tout $X$ dans $M$, on considère la
\deleted{\unsure{a) l}} catégorie $C_{/X}\simeq C_{/\varphi(\unsure{x|X})}$, et la flèche
canonique dans \deleted{\unsure{la}} $M^\vee$
\[
\varinjlim_{x\in C_{/X}} i(x)\longrightarrow X\mathpunct{.}
\tag*{(4.8.2.2)}
\]

$\unsure{1|2}^{\mathrm{o}}$) Considérons les conditions suivantes

\textcircled{a} \inserted{Pour tt $X\in\operatorname{Ob}M$,} la famille des flèches $i(x)\longrightarrow X$, indexée
par $\operatorname{Ob}C_{/X}$, est épimorphique stricte.
